%! AP Statistics — Unit 7, Lesson 7.1 — Group Activity (Answer Key)
\documentclass[10pt]{article}
\usepackage{apstats-article}
\usepackage{apstats-key}

\begin{document}

\pageheader{Unit 7, Lesson 7.1}{Group Activity --- Answer Key}
\namepartnerperiod

\vspace{0.1in}
\textbf{Directions:} Work with your partner on each scenario. For each, (a) state the
hypotheses, (b) describe a Type I error in context, (c) describe a Type II error in
context, and (d) decide which error is more costly and explain why.

\begin{tcolorbox}[colback=white, colframe=black!40, title=\textbf{Tier R --- Remediate},
    fonttitle=\bfseries, arc=2mm, left=3mm, right=3mm, top=2mm, bottom=2mm]

\textbf{Scenario 1 --- School Nurse Study}\\
A school nurse believes the mean number of sick days per year for students at her school
differs from the district average of 5 days.

\begin{enumerate}[label=\alph*.]
  \item State the hypotheses:
        $H_0$: The mean number of sick days is \ans{5} days.
        $H_a$: The mean number of sick days is \ans{not equal to} 5 days.

  \item Describe a Type I error in context:
        \ansline{Concluding the mean number of sick days differs from 5 when it actually equals 5 days.}

  \item Describe a Type II error in context:
        \ansline{Failing to conclude the mean differs from 5 days when it actually does differ.}

  \item Which error is more costly?
        \begin{teachernote}
            Both are defensible. Type I: the nurse wastes resources investigating a
            non-existent problem. Type II: a real health concern goes unaddressed.
        \end{teachernote}
        \ansline{Answers vary --- accept either with reasonable justification.}
\end{enumerate}
\end{tcolorbox}

\begin{tcolorbox}[colback=white, colframe=black!40, title=\textbf{Tier A --- Approaching Proficiency},
    fonttitle=\bfseries, arc=2mm, left=3mm, right=3mm, top=2mm, bottom=2mm]

\textbf{Scenario 2 --- Bottling Plant Quality Control}\\
A factory bottles fruit juice. The label states 12 fl oz per bottle.

\begin{enumerate}[label=\alph*.]
  \item State the hypotheses:
        \ansline{$H_0{:}\ \mu = 12$ fl oz; $H_a{:}\ \mu \ne 12$ fl oz}

  \item Describe a Type I error in context:
        \ansline{Concluding the mean fill differs from 12 fl oz when it actually equals 12 fl oz --- the factory halts production unnecessarily.}

  \item Describe a Type II error in context:
        \ansline{Failing to conclude the mean fill differs from 12 fl oz when it actually does --- underfilled or overfilled bottles continue to be sold.}

  \item Which error is more costly?
        \ansline{For the company: Type I (unnecessary production shutdowns cost money). For consumers: Type II (they receive systematically underfilled bottles). Both perspectives are valid.}
\end{enumerate}
\end{tcolorbox}

\begin{tcolorbox}[colback=white, colframe=black!40, title=\textbf{Tier E --- Extension},
    fonttitle=\bfseries, arc=2mm, left=3mm, right=3mm, top=2mm, bottom=2mm]

\textbf{Scenario 3 --- Clinical Drug Trial}

\begin{enumerate}[label=\alph*.]
  \item Error types in context:
        \ansline{Type I ($\alpha = 0.01$): approving an antibiotic that does not actually reduce infection duration below 7 days --- patients are given an ineffective drug. Type II: failing to approve an effective antibiotic --- patients miss out on a beneficial treatment.}

  \item Counterargument to raising $\alpha$ to 0.10:
        \ansline{Raising $\alpha$ to 0.10 increases the Type I error rate tenfold --- a 10\% chance of approving an ineffective drug. In a clinical context, the consequences of treating patients with a drug that doesn't work (and potentially has side effects) can be severe. Keeping $\alpha = 0.01$ prioritizes patient safety.}

  \item Effect of increasing $n$:
        \ansline{With $\alpha$ held constant at 0.01 and a larger sample ($n = 200$), the test has more power to detect a real effect. The Type II error rate $\beta$ decreases and power $= 1 - \beta$ increases. The Type I error rate remains 0.01 (controlled by $\alpha$, not $n$).}
\end{enumerate}
\end{tcolorbox}

\end{document}
